\section{Posterior coordinates and two local inequalities}

An experiment observes a random variable about a uniform input spin.
Its posterior law is the law of its posterior vector $p$ under the marginal
observation distribution. It is symmetric if it is invariant under
permutations of the four spin labels. Set
\[
L=\begin{pmatrix}
1&1&-1&-1\\
1&-1&1&-1\\
1&-1&-1&1
\end{pmatrix},\qquad x=Lp,
\]
and define
\begin{equation}\label{eq:statistics}
m(\mu)=\E_\mu|x|^2,\qquad
\tau(\mu)=\E_\mu x_1x_2x_3,\qquad h(\mu)=|\tau(\mu)|.
\end{equation}
The following bounds will be used repeatedly:
\begin{equation}\label{eq:bounds}
|x|^2=4\|p-\uvec\|_2^2,\quad 0\le m\le3,\quad h\le \frac m3,
\qquad \frac m8\le\E\TV(p,\uvec)\le\frac12\sqrt m.
\end{equation}
The cubic bound follows pointwise from
$3|x_1x_2x_3|\le x_1^2+x_2^2+x_3^2$, by AM--GM and $|x_i|\le1$.
The lower total-variation bound follows from
$\sum_i(p_i-1/4)^2\le\sum_i|p_i-1/4|$, while the upper bound follows from
Cauchy--Schwarz and Jensen. If an observation is degraded, its posterior
becomes a conditional expectation of the old posterior, so $m$ cannot
increase.

A channel with parameter $t\in[0,1]$ changes the posterior to
$\uvec+t(p-\uvec)$ and hence
\begin{equation}\label{eq:scaling}
x\mapsto tx,\qquad m\mapsto t^2m,\qquad
\tau\mapsto t^3\tau,\qquad h\mapsto t^3h.
\end{equation}
The revealing experiment has $m=3$ and $\tau=1$.

\begin{proposition}[Imported moment-saving input]\label{prop:input}
There is a class $\calC$ of symmetric posterior laws containing the
revealing and uninformative laws and closed under every channel step
$t\in[0,1]$ and under conditionally independent products. If
$\mu,\nu\in\calC$, write $u=m(\mu)$ and $v=m(\nu)$, and let $x=Lp$ and
$y=Lr$ under their respective posterior laws. Then
\begin{equation}\label{eq:input-strong}
 m(\mu\star\nu)
 \le u+v-\frac{\E_\mu|x|^4\,v+u\,\E_\nu|y|^4}{1000}.
\end{equation}
Consequently, with $c=1/1000$,
\begin{equation}\label{eq:input}
 m(\mu\star\nu)\le u+v-cuv(u+v).
\end{equation}
\end{proposition}

The first assertion is Proposition~3.2 and its closure statement in
\cite{OAIThreshold}. It explicitly permits arbitrary channel parameters and
nonidentical input laws. The consequence \eqref{eq:input} follows from
Jensen's inequalities $\E|x|^4\ge u^2$ and $\E|y|^4\ge v^2$. Every finite
subtree broadcast experiment and every partial product of its branches
belongs to $\calC$. Limiting posterior laws will not need to be in $\calC$.

The next lemma, unlike Proposition~\ref{prop:input}, needs only symmetry.
It strengthens the smoothing-dependent estimate in
\cite[Lemma~4.1]{OAICapacity}.

\begin{lemma}[Universal local moment bounds]\label{lem:local}
For any two symmetric posterior laws, put $u=m(\mu)$ and $v=m(\nu)$. Then
\begin{align}
 m(\mu\star\nu)
 &\ge u+v-5\bigl(h(\mu)v+u h(\nu)\bigr),\label{eq:mlower}\\
 h(\mu\star\nu)
 &\le h(\mu)+h(\nu)+376uv.\label{eq:hupper}
\end{align}
No smoothing or membership in $\calC$ is assumed.
\end{lemma}

\begin{proof}
Take independent marginal posterior vectors $p,r$, and write
$x=Lp$, $y=Lr$, and $s=x\cdot y$. The joint marginal distribution of the
observations is their independent marginal product weighted by
(cf. the posterior-product formalism in \cite{MM06})
\[
 Z=4\sum_i p_i r_i=1+s.
\]
When $Z>0$, the combined posterior has coordinates
\begin{equation}\label{eq:bayes}
 LP=\frac NZ,\qquad N=x+y+R,\qquad
 R_i=x_jy_k+x_ky_j,\quad\{i,j,k\}=\{1,2,3\}.
\end{equation}
In particular, $|N_i|\le Z$ and $0\le Z\le4$. At $Z=0$, all $N_i$
vanish; all weighted posterior quantities below are assigned zero. Their
bounds show that these extensions are continuous.

Symmetry acts on $x$ by coordinate permutations and even sign changes.
Consequently
\[
 \E x_i=0,\qquad \E x_ix_j=\frac u3\mathbf1_{\{i=j\}},
\]
and among degree-three monomials only $x_1x_2x_3$ can have nonzero
expectation. Put $\zeta_\mu=\tau(\mu)$ and $\zeta_\nu=\tau(\nu)$.
Direct expansion gives
\[
 \E|R|^2=\frac23uv,\qquad \E s^2=\frac13uv,
\]
\[
 \E[s\,x\cdot R]=2\zeta_\mu v,\qquad
 \E[s\,y\cdot R]=2u\zeta_\nu,\qquad
 \E[s|R|^2]=6\zeta_\mu\zeta_\nu.
\]
Using $Z^{-1}=1-s+s^2/Z$, with the marginal weight $Z$ included, we
obtain the exact identity
\begin{equation}\label{eq:exactm}
 m(\mu\star\nu)
 =u+v-4(\zeta_\mu v+u\zeta_\nu)-6\zeta_\mu\zeta_\nu
 +\E\frac{|N|^2s^2}{Z}.
\end{equation}
The mixed quadratic terms cancel since $\E|R|^2=2\E s^2$, and the remainder
is nonnegative. By $h(\mu)\le u/3$ and $h(\nu)\le v/3$,
\[
 6|\zeta_\mu\zeta_\nu|
 \le h(\mu)v+u h(\nu).
\]
Equation \eqref{eq:mlower} follows.

For the third moment,
\[
 \tau(\mu\star\nu)=\E\frac{N_1N_2N_3}{Z^2},\qquad
 Z^{-2}-(1-2s)=\frac{s^2(3+2s)}{Z^2}.
\]
Because $|N_i|\le Z$, $-1\le s\le3$, and $Z\le4$, the absolute remainder
after replacing $Z^{-2}$ by $1-2s$ is at most
\[
 36\E s^2=12uv.
\]
An exact expansion of $N_1N_2N_3(1-2s)$ has two pure terms,
$x_1x_2x_3$ and $y_1y_2y_3$, each with coefficient one. Mixed terms of
degree one in either vector average to zero. The sum of the absolute
coefficients of all remaining mixed terms is $364$. Every such term has
degree at least two in each vector and absolute expectation at most $uv$:
all coordinate absolute values are at most one, and retaining two factors
bounds each marginal monomial by the corresponding squared norm. Hence the
polynomial expectation differs from $\zeta_\mu+\zeta_\nu$ by at most
$364uv$. Adding the remainder proves \eqref{eq:hupper}. The coefficient
count is an exact integer expansion and is reproduced by the verification
script accompanying the source.

All expressions at $Z=0$ extend as asserted because they are bounded by a
constant times $Z$, possibly multiplied by a bounded power of $s$.
\end{proof}
